
\documentclass[12pt,a4paper]{article}

\usepackage{fontspec}
\usepackage{xeCJK}
\usepackage{polyglossia}
\setmainlanguage{russian}
\setotherlanguage{english}
\setmainfont{DejaVu Serif}
\setsansfont{DejaVu Sans}
\setmonofont{DejaVu Sans Mono}
\setCJKmainfont{Noto Serif CJK SC}

\usepackage{amsmath,amssymb,amsthm,mathtools}
\usepackage{geometry}
\geometry{left=18mm,right=18mm,top=18mm,bottom=20mm}
\usepackage{enumitem}
\setlist{itemsep=0.15em,topsep=0.25em,leftmargin=1.4em}
\usepackage{xcolor}
\usepackage{array,booktabs,longtable}
\usepackage{microtype}
\usepackage{titlesec}
\usepackage{fancyhdr}
\usepackage{hyperref}
\hypersetup{colorlinks=true,linkcolor=black,urlcolor=blue}

\pagestyle{fancy}
\setlength{\headheight}{15pt}
\fancyhf{}
\lhead{Государственный экзамен 2026}
\rhead{Краткие ответы, кафедра ОУ ВМК МГУ}
\cfoot{\thepage}

\setlength{\parindent}{0pt}
\setlength{\parskip}{0.38em}
\allowdisplaybreaks
\sloppy
\setcounter{secnumdepth}{0}
\setcounter{tocdepth}{2}
\titleformat{\section}{\Large\bfseries\raggedright}{\thesection.}{0.5em}{}
\titleformat{\subsection}{\large\bfseries\raggedright}{\thesubsection.}{0.5em}{}
\titleformat{\subsubsection}{\normalsize\bfseries\raggedright}{\thesubsubsection.}{0.5em}{}

\newcommand{\R}{\mathbb{R}}
\newcommand{\eps}{\varepsilon}
\newcommand{\T}{\top}
\newcommand{\dd}{\,d}
\newcommand{\norm}[1]{\left\lVert #1\right\rVert}
\newcommand{\inner}[2]{\left\langle #1,#2\right\rangle}
\newcommand{\Lie}{\operatorname{Lie}}
\newcommand{\Vect}{\operatorname{Vec}}
\newcommand{\rank}{\operatorname{rank}}
\newcommand{\ad}{\operatorname{ad}}
\newcommand{\sgn}{\operatorname{sgn}}
\newcommand{\sng}{\mathrm{sng}}
\newcommand{\PV}{\mathrm{PV}}
\newcommand{\CV}{\mathrm{CV}}
\DeclareMathOperator*{\argmax}{arg\,max}
\DeclareMathOperator*{\argmin}{arg\,min}
\newcommand{\Argmax}{\operatorname*{Arg\,max}}
\newcommand{\ticket}[2]{\section*{Билет #1. #2}\addcontentsline{toc}{section}{Билет #1. #2}}
\newcommand{\official}[1]{\subsection*{Официальная формулировка}\textbf{#1}}
\newcommand{\blocktitle}[3]{%
  \clearpage
  \thispagestyle{fancy}%
  \begin{center}
    {\LARGE\bfseries Блок #1. #2}\\[0.35em]
    {\large #3}
  \end{center}
  \addcontentsline{toc}{section}{Блок #1. #2}
  \vspace{0.45em}
}

\begin{document}
\begin{titlepage}
\centering
\vspace*{3cm}
{\Huge\bfseries Краткие ответы к государственному экзамену\\[0.4em]кафедры оптимального управления ВМК МГУ}\\[1.2cm]
{\Large Билеты 1--29}\\[0.8cm]
{\large Объединенная версия по блокам A--I}\\[2cm]
{\large 2026}
\vfill
\end{titlepage}

\tableofcontents

\blocktitle{A}{Обратные задачи и метод динамической регуляризации}{Билеты 1--3. Краткая версия для запоминания}

\section*{Билет 1. Обратные задачи для систем ОДУ. Метод динамической регуляризации}

Рассматривается система
\[
\dot y(t)=f(t,y(t))u(t)+g(t,y(t)),\qquad y(0)=y_0,\qquad 0\le t\le T,
\]
где \(y(t)\in\mathbb R^n\), \(u(t)\in P\subset\mathbb R^r\),
\[
U=\{u(\cdot)\in L_2^r(0,T):u(t)\in P\ \text{п.в.}\}.
\]

\textbf{Прямая задача:} по заданному \(u(\cdot)\) найти \(y(\cdot)\).  
\textbf{Обратная задача:} по наблюдаемой траектории \(y(\cdot)\) восстановить \(u(\cdot)\).

Обратная задача некорректна: управление может быть неединственным и неустойчивым относительно ошибок наблюдения. Поэтому вводится \textbf{нормальное управление}:
\[
u_*=\operatorname*{arg\,min}\{\|u\|_{L_2^r(0,T)}:\ y(t,u)=y^r(t),\ 0\le t\le T\}.
\]

На практике известны дискретные наблюдения:
\[
|\tilde y_i-y^r(t_i)|\le \delta,\qquad
0=t_0<t_1<\cdots<t_N=T,\qquad
h=\max_i(t_{i+1}-t_i).
\]

Метод динамической регуляризации строит кусочно-постоянное управление
\[
\tilde u_{\delta,h}(t)=u_i,\qquad t\in[t_i,t_{i+1}),\qquad u_i\in P.
\]

Вводится траектория-поводырь:
\[
z_0=\tilde y_0,\qquad
\frac{z_{i+1}-z_i}{t_{i+1}-t_i}
=f(t_i,\tilde y_i)u_i+g(t_i,\tilde y_i).
\]

На каждом шаге \(u_i\) выбирается из регуляризованной задачи:
\[
T_i(v)=
2\left\langle z_i-\tilde y_i,\,
f(t_i,\tilde y_i)v+g(t_i,\tilde y_i)
\right\rangle+\alpha |v|^2
\to \min_{v\in P},
\]
\[
T_i(u_i)\le \inf_{v\in P}T_i(v)+\varepsilon.
\]

Член \(\alpha|v|^2\) --- стабилизатор Тихонова; он подавляет неустойчивость и обеспечивает стремление к управлению минимальной нормы.

При согласованном стремлении
\[
\alpha,\delta,h,\varepsilon\to0,\qquad
\frac{\delta+h+\varepsilon}{\alpha}\to0,
\]
имеем
\[
z_h\to y^r \quad \text{в } C[0,T],
\qquad
\tilde u_{\delta,h}\to u_* \quad \text{в } L_2^r(0,T).
\]

\textbf{Ключевая фраза:} метод динамической регуляризации заменяет некорректное восстановление управления устойчивой пошаговой процедурой с поводырём и стабилизатором Тихонова.

\section*{Билет 2. Обратные задачи для параболических уравнений}

Типичный пример --- уравнение теплопроводности:
\[
\begin{cases}
y_t(t,x)=a^2y_{xx}(t,x)+u(t,x), &(t,x)\in(0,T)\times(0,l),\\
y(t,0)=y(t,l)=0,\\
y(0,x)=y_0(x).
\end{cases}
\]

\textbf{Прямая задача:} по \(u(t,x)\) найти \(y(t,x)\).  
\textbf{Обратная задача:} по наблюдениям \(y(t,x)\) восстановить правую часть \(u(t,x)\).

Используется гильбертов тройник:
\[
V=H_0^1(0,l),\qquad H=L_2(0,l),\qquad
V\subset H\simeq H^*\subset V^*.
\]

Оператор
\[
A:V\to V^*
\]
задаётся слабой формой:
\[
\langle Ay,v\rangle_{V^*,V}
=
a^2\int_0^l y_xv_x\,dx.
\]

Тогда уравнение записывается как
\[
\dot y(t)+Ay(t)=u(t),\qquad y(0)=y_0.
\]

Слабое решение:
\[
y\in L_2(0,T;V),\qquad \dot y\in L_2(0,T;V^*),
\qquad y\in C([0,T];H).
\]

Допустимые управления:
\[
U=\{u\in L_2(0,T;H):u(t)\in P\ \text{п.в.}\},
\]
где \(P\subset H\) выпукло, замкнуто и ограничено.

Нормальное управление:
\[
u_*=\operatorname*{arg\,min}\{\|u\|_{L_2(0,T;H)}:\ y(t,u)=y^r(t)\}.
\]

Наблюдения:
\[
\|y^r(t_i)-\tilde y_i\|_H\le\delta,\qquad
h=\max_i(t_{i+1}-t_i).
\]

Метод строит кусочно-постоянное управление:
\[
\tilde u_{\delta,h}(t)=u_i,\qquad t\in[t_i,t_{i+1}),\qquad u_i\in P.
\]

Поводырь:
\[
\dot z_h(t)+Az_h(t)=u_i,\qquad t\in(t_i,t_{i+1}).
\]

На шаге \(i\):
\[
\varphi_i(v)=
2\langle z_h(t_i)-\tilde y_i,v\rangle_H+\alpha\|v\|_H^2
\to\min_{v\in P},
\]
\[
\varphi_i(u_i)\le \inf_{v\in P}\varphi_i(v)+\varepsilon.
\]

При согласованном стремлении параметров к нулю:
\[
z_h\to y^r \quad \text{в } C([0,T];H),
\qquad
\tilde u_{\delta,h}\to u_* \quad \text{в } L_2(0,T;H).
\]

\textbf{Важно:} в параболическом случае нельзя смешивать \(L_2\)-скалярное произведение и двойственность \(V^*,V\). Оператор \(A\) корректно задаётся как \(A:V\to V^*\).

\textbf{Замечание:} точная форма оценок по \(h,\delta,\varepsilon,\alpha\) требует сверки с оригинальным текстом.

\section*{Билет 3. Обратные задачи для гиперболических уравнений}

Типичный пример --- волновое уравнение:
\[
\begin{cases}
y_{tt}(t,x)=a^2y_{xx}(t,x)+u(t,x), &(t,x)\in(0,T)\times(0,l),\\
y(t,0)=y(t,l)=0,\\
y(0,x)=y_0(x),\\
y_t(0,x)=y_1(x).
\end{cases}
\]

\textbf{Прямая задача:} по \(u(t,x)\) найти \(y(t,x)\).  
\textbf{Обратная задача:} по наблюдениям \(y\) или \(y_t\) восстановить \(u(t,x)\).

Главное отличие от параболического случая: уравнение второго порядка по времени, поэтому состояние есть пара
\[
X(t)=(y(t),\dot y(t)).
\]

Пространства:
\[
V=H_0^1(0,l),\qquad H=L_2(0,l),\qquad
V\subset H\simeq H^*\subset V^*.
\]

Оператор:
\[
\langle Ay,v\rangle_{V^*,V}
=
a^2\int_0^l y_xv_x\,dx.
\]

Абстрактная форма:
\[
\ddot y(t)+Ay(t)=u(t),
\qquad
y(0)=y_0,\quad \dot y(0)=y_1.
\]

Энергетический класс решений:
\[
y\in C([0,T];V),\qquad
\dot y\in C([0,T];H).
\]

Нормальное управление:
\[
u_*=\operatorname*{arg\,min}\{\|u\|_{L_2(0,T;H)}:\ y(t,u)=y^r(t)\}.
\]

Если наблюдается скорость, то данные имеют вид
\[
\|\dot y^r(t_i)-\tilde w_i\|_H\le\delta.
\]
Если наблюдается состояние, то
\[
\|y^r(t_i)-\tilde y_i\|_H\le\delta.
\]
В ответе важно не смешивать эти два типа наблюдений.

Метод динамической регуляризации строит
\[
\tilde u_{\delta,h}(t)=u_i,\qquad t\in[t_i,t_{i+1}),\qquad u_i\in P.
\]

Поводырь удовлетворяет второму порядку:
\[
\ddot z_h(t)+Az_h(t)=u_i,\qquad t\in(t_i,t_{i+1}).
\]

Если наблюдается скорость, то типичная регуляризованная задача:
\[
\varphi_i(v)
=
2\langle \dot z_h(t_i)-\tilde w_i,v\rangle_H
+
\alpha\|v\|_H^2
\to\min_{v\in P},
\]
\[
\varphi_i(u_i)\le \inf_{v\in P}\varphi_i(v)+\varepsilon.
\]

Стабилизатор Тихонова \(\alpha\|v\|_H^2\) обеспечивает устойчивость и выбор нормального решения.

При согласованном стремлении параметров к нулю:
\[
(z_h,\dot z_h)\to (y^r,\dot y^r)
\]
в соответствующей энергетической норме, а
\[
\tilde u_{\delta,h}\to u_*
\quad\text{в }L_2(0,T;H).
\]

\textbf{Ключевое отличие:} для гиперболического уравнения сходимость надо формулировать не только для \(y\), а для пары \((y,\dot y)\) в энергетическом пространстве.

\textbf{Замечание:} точная форма алгоритма и оценок зависит от того, что именно наблюдается: \(y\), \(\dot y\) или оба компонента. Это требует сверки с оригинальным источником.

\section*{Общая схема блока A}

\[
\begin{gathered}
\text{некорректная обратная задача}\Rightarrow\text{нормальное управление}\\
\Rightarrow\text{дискретные наблюдения}\Rightarrow\text{поводырь}\\
\Rightarrow\text{регуляризованная минимизация}\\
\Rightarrow\text{сходимость к нормальному управлению}.
\end{gathered}
\]


\blocktitle{B}{Экономический рост и оптимальное управление}{Билеты 4--8. Краткая версия для запоминания}

\section*{Билет 4. Одномерная модель Рамсея. Cobb--Douglas. Особый режим}

Модель: капитал $x(t)>0$, управление $u(t)\in[0,1]$ -- доля выпуска, идущая в инвестиции. Доля $1-u(t)$ идет в потребление.
\[
\dot x=uF(x)-\mu x,\qquad x(0)=x_0>0,
\]
\[
F(x)=Ax^\varepsilon,\qquad A>0,\quad 0<\varepsilon<1.
\]
Функционал на бесконечном горизонте:
\[
J[u]=\int_0^\infty e^{-\nu t}(1-u(t))F(x(t))\,dt\to\max .
\]

Гамильтониан:
\[
H=e^{-\nu t}(1-u)F(x)+\psi\bigl(uF(x)-\mu x\bigr).
\]
Коэффициент при $u$:
\[
F(x)(\psi-e^{-\nu t}).
\]
Поэтому функция переключения:
\[
\Phi(t)=\psi(t)-e^{-\nu t}.
\]
\[
\Phi>0\Rightarrow u=1,\qquad \Phi<0\Rightarrow u=0.
\]

В текущих переменных
\[
p(t)=e^{\nu t}\psi(t),
\]
текущий гамильтониан:
\[
M=(1-u)F(x)+p(uF(x)-\mu x).
\]
Особый режим:
\[
p\equiv1.
\]
Из уравнения $\dot p=\nu p-M_x$ при $p=1$ получаем
\[
F'(x_s)=\mu+\nu.
\]
Особое управление определяется из $\dot x=0$:
\[
u_sF(x_s)=\mu x_s,
\qquad
u_s=\frac{\mu x_s}{F(x_s)}.
\]
Для $F(x)=Ax^\varepsilon$:
\[
x_s=\left(\frac{A\varepsilon}{\mu+\nu}\right)^{\frac{1}{1-\varepsilon}},
\qquad
u_s=\frac{\mu\varepsilon}{\mu+\nu}.
\]
В частном случае $A=1$, $\varepsilon=\frac12$:
\[
x_s=\frac{1}{4(\mu+\nu)^2},
\qquad
u_s=\frac{\mu}{2(\mu+\nu)}.
\]

Качественно:
\[
x_0<x_s\Rightarrow u=1\ \text{до выхода на }x_s,
\]
\[
x_0>x_s\Rightarrow u=0\ \text{до выхода на }x_s,
\]
\[
x_0=x_s\Rightarrow x(t)\equiv x_s,
\quad u(t)\equiv u_s.
\]
Смысл: оптимальная политика выводит экономику на уровень, где
\[
F'(x_s)=\mu+\nu,
\]
то есть предельная производительность капитала равна выбытию плюс дисконтирование.

\section*{Билет 5. Двухсекторная модель Cobb--Douglas. Особый режим}

Модель распределения ресурсов:
\[
\dot x_1=uF(x),\qquad \dot x_2=(1-u)F(x),
\qquad 0\le u\le1,
\]
\[
x_1(0)=x_{10}>0,
\qquad x_2(0)=x_{20}>0,
\]
\[
J[u]=x_2(T)\to\max.
\]
Производственная функция Cobb--Douglas:
\[
F(x)=Ax_1^{\varepsilon_1}x_2^{\varepsilon_2},
\qquad
\varepsilon_1>0,
\quad \varepsilon_2>0,
\quad \varepsilon_1+\varepsilon_2=1.
\]
Множитель $A>0$ можно убрать заменой масштаба времени.

Гамильтониан:
\[
H=\psi_1uF(x)+\psi_2(1-u)F(x)
=F(x)\bigl(u(\psi_1-\psi_2)+\psi_2\bigr).
\]
Терминальные условия:
\[
\psi_1(T)=0,
\qquad
\psi_2(T)=1.
\]
Функция переключения:
\[
\Pi=\psi_1-\psi_2.
\]
\[
\Pi>0\Rightarrow u=1,
\qquad
\Pi<0\Rightarrow u=0,
\qquad
\Pi=0\Rightarrow \text{особый режим}.
\]

Сопряженная система:
\[
\dot\psi_1=-F_{x_1}'(x)(u\Pi+\psi_2),
\qquad
\dot\psi_2=-F_{x_2}'(x)(u\Pi+\psi_2).
\]
На особой дуге $\Pi\equiv0$, следовательно $\psi_1=\psi_2$ и $\dot\psi_1=\dot\psi_2$. Поэтому
\[
F_{x_1}'(x)=F_{x_2}'(x).
\]
Для Cobb--Douglas:
\[
F_{x_1}'(x)=\frac{\varepsilon_1}{x_1}F(x),
\qquad
F_{x_2}'(x)=\frac{\varepsilon_2}{x_2}F(x).
\]
Значит,
\[
\frac{x_1}{\varepsilon_1}=\frac{x_2}{\varepsilon_2}.
\]
Особая магистраль:
\[
L_{\rm sng}=\left\{x\in\mathbb R_+^2:\frac{x_1}{\varepsilon_1}=\frac{x_2}{\varepsilon_2}\right\}.
\]
Чтобы оставаться на ней,
\[
\frac{\dot x_1}{\varepsilon_1}=\frac{\dot x_2}{\varepsilon_2}.
\]
Отсюда
\[
\frac{u}{\varepsilon_1}=\frac{1-u}{\varepsilon_2}
\quad\Rightarrow\quad
u_{\rm sng}=\varepsilon_1.
\]

Удобная замена координат:
\[
y_1=\frac{x_1}{\varepsilon_1},
\qquad
y_2=\frac{x_2}{\varepsilon_2}.
\]
В этих координатах
\[
L_{\rm sng}:\quad y_1=y_2.
\]

Для достаточно большого $T$ структура оптимального управления:
\[
\text{граничный участок}\to \text{особая магистраль}\to \text{терминальный участок}.
\]
То есть
\[
u(t)=
\begin{cases}
0\ \text{или}\ 1, & 0\le t<\tau,\\
\varepsilon_1, & \tau<t<\theta,\\
0, & \theta<t\le T.
\end{cases}
\]
Если начальная точка уже на магистрали, то $\tau=0$.

Обоснование оптимальности: максимизированный гамильтониан
\[
M(x,\psi)=F(x)g(\psi),
\qquad
g(\psi)=\max\{\psi_1,\psi_2\}
=\frac{\psi_1+\psi_2+|\psi_1-\psi_2|}{2}.
\]
Так как $F$ вогнута и $g(\psi(t))>0$, для любой допустимой пары $(\widehat x,\widehat u)$:
\[
\Delta J\le
\int_0^T
\bigl[F(\widehat x)-F(x)-(F'(x),\widehat x-x)\bigr]g(\psi)\,dt\le0.
\]
Следовательно, экстремальный процесс оптимален.

Главные формулы для запоминания:
\[
L_{\rm sng}:\frac{x_1}{\varepsilon_1}=\frac{x_2}{\varepsilon_2},
\qquad
u_{\rm sng}=\varepsilon_1.
\]

\section*{Билет 6. ПМП на бесконечном интервале времени}

Типичная задача:
\[
\dot x=f(x,u),
\qquad
u(t)\in U,
\qquad
x(0)=x_0,
\]
\[
J[x,u]=\int_0^\infty e^{-\rho t}g(x(t),u(t))\,dt\to\max.
\]
Функционал является несобственным интегралом, поэтому нужны условия сходимости, например
\[
\int_T^\infty e^{-\rho t}|g(x(t),u(t))|\,dt\to0.
\]

Гамильтониан Понтрягина:
\[
H(x,t,u,\psi,\psi^0)=\langle f(x,u),\psi\rangle+
\psi^0e^{-\rho t}g(x,u),
\qquad \psi^0\ge0.
\]
Нетривиальность:
\[
\|\psi(0)\|+\psi^0>0.
\]

Если $(x^*,u^*)$ оптимальна, то существуют $\psi(t)$ и $\psi^0\ge0$ такие, что
\[
\dot x^*=f(x^*,u^*),
\]
\[
\dot\psi=-H_x(x^*,t,u^*,\psi,\psi^0),
\]
то есть
\[
\dot\psi= -f_x(x^*,u^*)^T\psi
-\psi^0 e^{-\rho t}g_x(x^*,u^*).
\]
Условие максимума:
\[
H(x^*,t,u^*,\psi,\psi^0)=\max_{u\in U}H(x^*,t,u,\psi,\psi^0)
\quad \text{п.в.}
\]

Если $\psi^0>0$, то нормальный случай, можно положить
\[
\psi^0=1.
\]
Если $\psi^0=0$, то анормальный случай; его нельзя исключать без дополнительных условий.

Главное отличие бесконечного горизонта: условия трансверсальности на бесконечности не следуют автоматически. Формулы
\[
\lim_{t\to\infty}\psi(t)=0,
\qquad
\lim_{t\to\infty}\langle \psi(t),x^*(t)\rangle=0
\]
не являются общими теоремами. Они требуют дополнительных предположений.

Запомнить:
\[
\begin{gathered}
\text{ПМП}=\text{сопряженная система}+\text{условие максимума}\\
+\text{нетривиальность}+\text{осторожная трансверсальность}.
\end{gathered}
\]

\section*{Билет 7. Беллман, ПМП и текущие сопряженные переменные}

Бесконечная стационарная задача:
\[
\dot x=f(x,u),
\qquad
J(x,u)=\int_0^\infty e^{-\rho t}g(x(t),u(t))\,dt\to\max.
\]
Функция Беллмана:
\[
V(x)=\sup_u\int_0^\infty e^{-\rho t}g(x(t),u(t))\,dt.
\]
Если $V\in C^1$, то уравнение Беллмана:
\[
\rho V(x)=\max_{u\in U}\{g(x,u)+\langle V_x(x),f(x,u)\rangle\}.
\]
Оптимальная обратная связь:
\[
u^*(x)\in\operatorname*{arg\,max}_{u\in U}
\{g(x,u)+\langle V_x(x),f(x,u)\rangle\}.
\]

Текущая сопряженная переменная:
\[
p(t)=V_x(x^*(t)).
\]
Связь с обычной present-value переменной:
\[
\psi(t)=e^{-\rho t}p(t),
\qquad
p(t)=e^{\rho t}\psi(t).
\]

Текущий гамильтониан:
\[
M(x,u,p)=g(x,u)+\langle p,f(x,u)\rangle.
\]
ПМП в текущих переменных:
\[
\dot x^*=f(x^*,u^*),
\]
\[
\dot p=\rho p-M_x(x^*,u^*,p),
\]
\[
M(x^*,u^*,p)=\max_{u\in U}M(x^*,u,p).
\]

Экономический смысл:
\begin{itemize}
\item $V(x)$ -- максимальная будущая дисконтированная полезность;
\item $V_x(x)$ -- предельная ценность состояния;
\item $p(t)$ -- текущая теневая цена;
\item $\psi(t)$ -- та же цена, приведенная к начальному моменту.
\end{itemize}

Важно:
\[
\text{Беллман дает достаточность через проверочную теорему,}
\]
\[
\text{а ПМП обычно дает только необходимые условия.}
\]

\section*{Билет 8. Достаточные условия оптимальности на бесконечном интервале}

Задача:
\[
\dot x=f(x,u),
\qquad
J=\int_0^\infty e^{-\rho t}g(x,u)\,dt\to\max.
\]
ПМП сам по себе обычно является только необходимым условием. Для достаточности нужны дополнительные условия.

Пусть допустимая пара $(x^*,u^*)$ удовлетворяет нормальной форме ПМП:
\[
\dot\psi=-H_x(x^*,t,u^*,\psi),
\]
\[
H(x^*,t,u^*,\psi)=\max_{u\in U}H(x^*,t,u,\psi),
\]
где
\[
H(x,t,u,\psi)=\langle f(x,u),\psi\rangle+e^{-\rho t}g(x,u).
\]

Максимизированный гамильтониан:
\[
\mathcal H(x,t,\psi(t))=
\max_{u\in U}\{\langle f(x,u),\psi(t)\rangle+e^{-\rho t}g(x,u)\}.
\]

Достаточные условия оптимальности:
\begin{enumerate}
\item область состояний $G$ выпукла;
\item функция $x\mapsto\mathcal H(x,t,\psi(t))$ вогнута по $x$;
\item выполнено условие на бесконечности:
\[
\liminf_{t\to\infty}\langle \psi(t),x(t)-x^*(t)\rangle\ge0
\]
для любой допустимой траектории $x(t)$.
\end{enumerate}
Тогда
\[
J[x,u]\le J[x^*,u^*],
\]
то есть $(x^*,u^*)$ оптимальна.

Идея доказательства:
\[
J_T[x,u]-J_T[x^*,u^*]
\le
-\langle \psi(T),x(T)-x^*(T)\rangle.
\]
Переходя к пределу $T\to\infty$ и используя предельное условие, получаем оптимальность.

Альтернатива -- проверочная теорема Беллмана: если существует гладкая функция $V$, удовлетворяющая
\[
\rho V(x)=\max_{u\in U}\{g(x,u)+\langle V_x(x),f(x,u)\rangle\},
\]
и управление достигает максимума, то оно оптимально при подходящем условии на бесконечности.

Запомнить:
\[
\text{ПМП}+\text{вогнутость максимизированного }H+\text{условие на }\infty
\Rightarrow \text{оптимальность}.
\]

\section*{Общий минимум для блока B}

\begin{itemize}
\item Билет 4: Ramsey: $F'(x_s)=\mu+\nu$, $u_sF(x_s)=\mu x_s$.
\item Билет 5: two-sector Cobb--Douglas: $x_1/\varepsilon_1=x_2/\varepsilon_2$, $u_{\rm sng}=\varepsilon_1$.
\item Билет 6: ПМП на бесконечности: сопряженное уравнение, максимум, нетривиальность; трансверсальность требует осторожности.
\item Билет 7: Bellman: $\rho V=\max\{g+\langle V_x,f\rangle\}$; $p=V_x$, $\psi=e^{-\rho t}p$.
\item Билет 8: достаточность: вогнутость максимизированного гамильтониана и условие на бесконечности.
\end{itemize}


\blocktitle{C}{Геометрическая теория управления}{Билеты 9--11. Краткая версия для запоминания}

\section*{Билет 9. Геометрическая теория управления}

Гладкое многообразие $M$ локально устроено как $\mathbb R^n$. В точке $q\in M$ есть касательное пространство $T_qM$.

Векторное поле:
\[
X:M\to TM,
\qquad X(q)\in T_qM.
\]
Оно задает ОДУ:
\[
\dot q=X(q)
\]
и поток $P_X^t$.

Управляемая система:
\[
\dot q=V_u(q),\qquad u\in U.
\]
Система, линейная по управлению:
\[
\dot q=f_0(q)+\sum_{i=1}^m u_i f_i(q).
\]
Бездрейфовая система:
\[
\dot q=\sum_{i=1}^m u_i f_i(q).
\]

Скобка Ли:
\[
[X,Y](\varphi)=X(Y\varphi)-Y(X\varphi).
\]
В координатах:
\[
[X,Y]^i=\sum_{j=1}^n
\left(
X^j\frac{\partial Y^i}{\partial x^j}
-
Y^j\frac{\partial X^i}{\partial x^j}
\right).
\]
Свойства:
\[
[X,Y]=-[Y,X],\qquad [X,X]=0,
\]
\[
[X,[Y,Z]]+[Y,[Z,X]]+[Z,[X,Y]]=0.
\]

Смысл: коммутаторы дают направления, возникающие из быстрых переключений управлений.

Алгебра Ли:
\[
\Lie\{f_1,\ldots,f_m\}
\]
содержит поля и все их повторные скобки:
\[
f_i,\quad [f_i,f_j],\quad [f_i,[f_j,f_k]],\quad\ldots
\]

Условие скобочной порождаемости:
\[
\Lie_q\{f_1,\ldots,f_m\}=T_qM
\qquad \forall q\in M.
\]
Если $\dim M=n$, то
\[
\rank\{f_i(q),[f_i,f_j](q),[f_i,[f_j,f_k]](q),\ldots\}=n.
\]

Для бездрейфовых систем это условие ведет к управляемости по теореме Чоу--Рашевского. Для систем с дрейфом
\[
\dot q=f_0(q)+\sum u_i f_i(q)
\]
ранговое условие обычно дает доступность, но не автоматически полную управляемость.

\textbf{Главная формула:}
\[
\dot q=\sum u_i f_i(q),
\qquad
\Lie_q\{f_1,\ldots,f_m\}=T_qM.
\]

\section*{Билет 10. Сопряженная точка и слабый минимум}

Рассматривается задача:
\[
J[x]=\int_{t_0}^{t_1} f(t,x,\dot x)\,dt\to\min,
\qquad x(t_0)=a,
\quad x(t_1)=b.
\]

Экстремаль удовлетворяет уравнению Эйлера--Лагранжа:
\[
\frac{d}{dt}f_{\dot x}(t,x,\dot x)-f_x(t,x,\dot x)=0.
\]

Для вариаций $x+\lambda h$, где
\[
h(t_0)=h(t_1)=0,
\]
вторая вариация:
\[
\delta^2J[h]=
\int_{t_0}^{t_1}
\left(
\langle A\dot h,\dot h\rangle
+
2\langle C\dot h,h\rangle
+
\langle Bh,h\rangle
\right)dt,
\]
где
\[
A=f_{\dot x\dot x},\qquad B=f_{xx},\qquad C=f_{x\dot x}.
\]

Условие Лежандра:
\[
A(t)\ge0.
\]
Усиленное условие Лежандра:
\[
A(t)>0.
\]

Уравнение Якоби:
\[
-\frac{d}{dt}(A\dot h+C^Th)+C\dot h+Bh=0.
\]

Точка $\tau\in(t_0,t_1]$ называется сопряженной с $t_0$, если существует нетривиальное решение уравнения Якоби:
\[
h(t_0)=0,
\qquad
h(\tau)=0,
\qquad
h\not\equiv0.
\]

Достаточное условие слабого минимума: если экстремаль удовлетворяет уравнению Эйлера, усиленному условию Лежандра и на $(t_0,t_1]$ нет сопряженных с $t_0$ точек, то она дает строгий слабый минимум.

\textbf{Главная схема:}
\[
\text{Euler}+
\text{strong Legendre}+
\text{no conjugate points}
\Rightarrow
\text{strict weak minimum}.
\]

\section*{Билет 11. Особые траектории и особые управления}

Рассмотрим систему, линейную по управлению:
\[
\dot x=f_0(x)+u f_1(x),
\qquad u\in[a,b].
\]

Hamiltonian:
\[
H=h_0+u h_1,
\qquad h_i(x,\psi)=\langle\psi,f_i(x)\rangle.
\]

Функция переключения:
\[
\Phi(t)=H_u=h_1(x(t),\psi(t)).
\]

Обычные участки:
\[
\Phi>0\Rightarrow u=b,
\qquad
\Phi<0\Rightarrow u=a.
\]

Особый режим:
\[
\Phi(t)\equiv0
\]
на некотором интервале. Тогда условие максимума не определяет управление.

Метод построения:
\[
\Phi=0,
\quad \dot\Phi=0,
\quad \ddot\Phi=0,
\quad\ldots
\]
до тех пор, пока в одной из производных явно не появится $u$. Тогда из этого уравнения находят
\[
u_{\sng}.
\]

Порядок особого управления $q$:
\[
\frac{\partial}{\partial u}
\left(\frac{d}{dt}\right)^k H_u=0,
\quad k=0,\ldots,2q-1,
\]
\[
\frac{\partial}{\partial u}
\left(\frac{d}{dt}\right)^{2q}H_u\ne0.
\]

Условие Kelley при максимизационной конвенции:
\[
(-1)^q
\frac{\partial}{\partial u}
\left(\frac{d}{dt}\right)^{2q}H_u<0.
\]

Через скобки Пуассона:
\[
\dot h=\{H,h\}.
\]
Для первого порядка:
\[
\Phi=h_1=0,
\]
\[
\dot\Phi=\{h_0,h_1\}=0,
\]
\[
\ddot\Phi=
\{h_0,\{h_0,h_1\}\}
+u\{h_1,\{h_0,h_1\}\}=0.
\]
Отсюда
\[
u_{\sng}
=-
\frac{\{h_0,\{h_0,h_1\}\}}
{\{h_1,\{h_0,h_1\}\}}.
\]

После нахождения обязательно проверить:
\[
u_{\sng}\in[a,b].
\]

\textbf{Главная схема:}
\[
\Phi\equiv0
\Rightarrow
\Phi=\dot\Phi=\ddot\Phi=\cdots=0
\Rightarrow
u_{\sng}
\Rightarrow
\text{проверка допустимости и Kelley}.
\]


\blocktitle{D}{Волновое уравнение, граничное управление и наблюдение}{Билеты 12--14. Краткая версия для запоминания}

\section*{Билет 12. Граничное управление и наблюдение}

Волновое уравнение:
\[
 y_{tt}=y_{xx}-q(x)y,\quad 0<x<l.
\]

Граничное управление:
\[
 y(t,0)=u_0(t),\qquad y(t,l)=u_1(t).
\]

Пространство управлений:
\[
 H=L_2(0,T)\times L_2(0,T).
\]

Оператор управления:
\[
 A:H\to F,
 \qquad
 Au=(y(T),y_t(T)).
\]

Типичная слабая постановка:
\[
 F=L_2(0,l)\times H^{-1}(0,l),
\]
\[
 F^*=H_0^1(0,l)\times L_2(0,l).
\]

Сопряженная задача:
\[
 p_{tt}=p_{xx}-q(x)p,
 \qquad p(t,0)=p(t,l)=0.
\]

Наблюдение:
\[
 A^*v=(p_x(\cdot,0),-p_x(\cdot,l)).
\]

Двойственность:
\[
 \langle Au,v\rangle_{F,F^*}
 =
 \langle u,A^*v\rangle_H.
\]

Неравенство наблюдаемости:
\[
 \|A^*v\|_{H^*}^2\ge \mu\|v\|_{F^*}^2,
 \qquad \mu>0.
\]

Критическое время:
\[
 T^*=l\quad \text{для двухстороннего управления},
\]
\[
 T^*=2l\quad \text{для одностороннего управления}.
\]

Главная схема:
\[
 T>T^*
 \Rightarrow
 \text{наблюдаемость}
 \Rightarrow
 \text{точная управляемость}.
\]

\section*{Билет 13. Неравенство наблюдаемости и вариационный метод}

Абстрактная задача управления:
\[
 Au=f,
 \qquad
 A:H\to F.
\]

Сопряженный оператор:
\[
 A^*:F^*\to H^*.
\]

Точная управляемость:
\[
 R(A)=F.
\]

Критерий точной управляемости:
\[
 \|A^*v\|_{H^*}^2
 \ge
 \mu\|v\|_{F^*}^2,
 \qquad \mu>0.
\]

Важно:
\[
 N(A^*)=\{0\}
\]
само по себе означает только плотность $R(A)$, то есть приближенную управляемость, но не точную.

Нормальное решение:
\[
 u_*=
 \operatorname*{arg\,min}\{\|u\|_H: Au=f\}.
\]

Истокопредставимость:
\[
 u_*=J_HA^*v_*.
\]

Вариационный функционал:
\[
 I(v)=\|A^*v\|^2-2\langle f,v\rangle\to\min.
\]

После нахождения $v_*$:
\[
 u_*=J_HA^*v_*.
\]

Конечномерная версия:
\[
 I_N(v_N)=\|A_N^*v_N\|^2-2\langle f_N,v_N\rangle\to\min.
\]

Главное условие дискретизации:
\[
 \langle A_Nu_N,v_N\rangle
 =
 \langle u_N,A_N^*v_N\rangle.
\]

\section*{Билет 14. Однородная струна и конечномерная схема}

Однородная струна:
\[
 p_{tt}=p_{xx},
 \qquad p(t,0)=p(t,l)=0.
\]

Энергия:
\[
 E(t)=\int_0^l(p_t^2+p_x^2)\,dx.
\]

Энергия сохраняется:
\[
 E(t)=E(0)=E(T).
\]

Наблюдение:
\[
 A^*v=(p_x(\cdot,0),-p_x(\cdot,l)).
\]

Неравенство наблюдаемости:
\[
 \int_0^T\left(|p_x(t,0)|^2+|p_x(t,l)|^2\right)dt
 \ge
 \mu\left(\|v_0\|_{H_0^1}^2+\|v_1\|_{L_2}^2\right).
\]

Для двухстороннего наблюдения:
\[
 T>l.
\]

В студенческой нормировке:
\[
 \mu=\frac{2(T-l)}{l}>0.
\]

Взаимно сопряженные конечномерные операторы:
\[
 A_N:H_N\to F_N,
 \qquad
 A_N^*:F_N^*\to H_N^*.
\]

Дискретная двойственность:
\[
 \langle A_Nu_N,v_N\rangle
 =
 \langle u_N,A_N^*v_N\rangle.
\]

Матрично:
\[
 A_N^*=G_H^{-1}A_N^TG_F.
\]

Квадратичная задача:
\[
 I_N(c)=c^TG_Nc-2b_N^Tc\to\min.
\]

Если $G_N>0$, то
\[
 G_Nc=b_N.
\]

Восстановление управления:
\[
 u_N=J_{H_N}A_N^*v_N.
\]

Ключевая схема блока D:
\[
 \langle Au,v\rangle=\langle u,A^*v\rangle
 \Rightarrow
 \|A^*v\|^2\ge\mu\|v\|^2
 \Rightarrow
 u_*=J_HA^*v_*
 \Rightarrow
 I(v)=\|A^*v\|^2-2\langle f,v\rangle.
\]


\blocktitle{E}{Пространства Соболева и оптимальное управление теплом}{Билеты 15--17. Краткая версия для запоминания}

\section*{Билет 15. Обобщенная производная и пространства Соболева}

Пусть \(\Omega\subset\mathbb R^n\), \(u\in L_{1,\mathrm{loc}}(\Omega)\). Функция \(v_\alpha\in L_{1,\mathrm{loc}}(\Omega)\) называется обобщенной производной \(D^\alpha u\), если
\[
\int_\Omega uD^\alpha\varphi\,dx
=
(-1)^{|\alpha|}
\int_\Omega v_\alpha\varphi\,dx
\quad
\forall\varphi\in C_c^\infty(\Omega).
\]
Тогда пишут
\[
D^\alpha u=v_\alpha.
\]

В одномерном случае:
\[
\int_a^b f(x)\varphi'(x)\,dx
=
-
\int_a^b \omega(x)\varphi(x)\,dx,
\quad
\forall\varphi\in C_c^\infty(a,b).
\]

Свойства: обобщенная производная единственна почти всюду, линейна, а если классическая производная существует, то она совпадает с обобщенной.

Пространство Соболева:
\[
W_p^m(\Omega)=
\{u\in L_p(\Omega):D^\alpha u\in L_p(\Omega),\ |\alpha|\le m\}.
\]

Норма:
\[
\|u\|_{W_p^m}=
\left(
\sum_{|\alpha|\le m}\|D^\alpha u\|_{L_p}^p
\right)^{1/p}.
\]

При \(p=2\):
\[
H^m(\Omega)=W_2^m(\Omega).
\]

Пример:
\[
H^1(a,b)=\{u\in L_2(a,b):u'\in L_2(a,b)\},
\]
\[
\|u\|_{H^1}^2=\int_a^b(u^2+(u')^2)\,dx.
\]

\[
H_0^1(\Omega)=\overline{C_c^\infty(\Omega)}^{H^1}.
\]

В одномерном случае:
\[
H^1(a,b)\hookrightarrow C[a,b],
\]
\[
u(\beta)-u(\alpha)=\int_\alpha^\beta u'(x)\,dx.
\]

Соболевское вложение: если \(mp<n\), то
\[
W_p^m(\Omega)\hookrightarrow L_q(\Omega),
\quad
q\le \frac{np}{n-mp}.
\]
Если \(mp>n\), то
\[
W_p^m(\Omega)\hookrightarrow C(\overline\Omega)
\]
или в гёльдеровы пространства.

Для тепловых задач:
\[
H_0^1\subset L_2\subset H^{-1}.
\]
Если
\[
y\in L_2(0,T;H_0^1),
\quad
y_t\in L_2(0,T;H^{-1}),
\]
то
\[
y\in C([0,T];L_2).
\]

Ключевая схема:
\[
\text{слабая производная}
\Rightarrow
\text{пространства Соболева}
\Rightarrow
\text{слабые решения PDE}.
\]

\newpage

\section*{Билет 16. Оптимальный нагрев стержня}

Пусть
\[
Q=(0,T)\times(0,l).
\]

Задача:
\[
J(u)=\int_0^l (x(s,T;u)-b(s))^2\,ds\to\inf_{u\in U},
\]
\[
x_t=x_{ss}+u,
\quad (s,t)\in Q,
\]
\[
x(0,t)=x(l,t)=0,
\quad
x(s,0)=\varphi(s),
\]
\[
U=\{u\in L_2(Q):\|u\|_{L_2(Q)}\le R\}.
\]

Функциональные пространства:
\[
V=H_0^1(0,l),
\quad H=L_2(0,l),
\quad V^*=H^{-1}(0,l),
\]
\[
V\subset H\simeq H^*\subset V^*.
\]

Обобщенное решение:
\[
x\in L_2(0,T;H_0^1(0,l)),
\]
\[
x_t\in L_2(0,T;H^{-1}(0,l)).
\]

Слабая форма:
\[
\langle x_t(t),v\rangle_{H^{-1},H_0^1}
+
\int_0^l x_s(s,t)v_s(s)\,ds
=
\int_0^l u(s,t)v(s)\,ds
\]
для всех
\[
v\in H_0^1(0,l).
\]

Важное вложение:
\[
x\in C([0,T];L_2(0,l)).
\]
Поэтому \(x(T)\) корректно определено.

Существование: для каждого
\[
u\in L_2(Q),
\quad
\varphi\in L_2(0,l)
\]
существует обобщенное решение.

Единственность: если \(x_1,x_2\) -- два решения, то \(z=x_1-x_2\) удовлетворяет
\[
z_t=z_{ss},
\quad z(0)=0.
\]
Энергетическая оценка:
\[
\frac12\frac{d}{dt}\|z(t)\|_{L_2}^2+
\|z_s(t)\|_{L_2}^2=0.
\]
Следовательно,
\[
z=0.
\]

Теорема Вейерштрасса:
\[
U=\{u:\|u\|\le R\}
\]
выпукло, замкнуто и ограничено в \(L_2(Q)\), значит слабо компактно.

Берем минимизирующую последовательность:
\[
u_k\rightharpoonup u_*\in U.
\]

Оператор
\[
u\mapsto x(T;u)
\]
линеен и непрерывен, поэтому
\[
J(u)=\|x(T;u)-b\|_{L_2}^2
\]
слабо полунепрерывен снизу.

Значит,
\[
J(u_*)=\inf_{u\in U}J(u).
\]

Итог:
\[
\begin{gathered}
\text{слабая компактность }U+\text{слабая полунепрерывность }J\\
\Rightarrow\text{существует оптимальное управление}.
\end{gathered}
\]

\newpage

\section*{Билет 17. Градиент функционала и РМПГ}

Задача нагрева:
\[
J(u)=\int_0^l (x(s,T;u)-b(s))^2ds\to\inf,
\]
\[
x_t=x_{ss}+u,
\quad x(0,t)=x(l,t)=0,
\quad x(s,0)=\varphi(s),
\]
\[
U=\{u\in L_2(Q):\|u\|\le R\}.
\]

Вариация состояния:
\[
z_t=z_{ss}+v,
\quad z(0,t)=z(l,t)=0,
\quad z(s,0)=0.
\]

Первая вариация:
\[
\delta J(u;v)=
2\int_0^l (x(s,T)-b(s))z(s,T)\,ds.
\]

Сопряженная задача:
\[
-\psi_t=\psi_{ss},
\]
\[
\psi(0,t)=\psi(l,t)=0,
\]
\[
\psi(s,T)=2(x(s,T;u)-b(s)).
\]

Тогда:
\[
\delta J(u;v)=\iint_Q \psi(s,t)v(s,t)\,dsdt.
\]

Следовательно:
\[
J'(u)=\psi.
\]

Критерий оптимальности на выпуклом \(U\):
\[
u_*\in U^*
\Longleftrightarrow
\langle J'(u_*),u-u_*\rangle\ge0
\quad \forall u\in U.
\]

Для шара
\[
U=\{u:\|u\|\le R\}
\]
эквивалентно:
\[
J'(u_*)+2\lambda_*u_*=0,
\]
\[
\lambda_*(\|u_*\|^2-R^2)=0,
\]
\[
\lambda_*\ge0,
\quad
\|u_*\|\le R.
\]

Если пишут
\[
J'(u_*)+\mu_*u_*=0,
\]
то
\[
\mu_*=2\lambda_*.
\]

Обычный метод проекции градиента:
\[
u_{k+1}=P_U(u_k-\beta_kJ'(u_k)).
\]

Функция Тихонова:
\[
T_k(u)=J(u)+\alpha_k\|u\|^2.
\]

Градиент:
\[
T_k'(u)=J'(u)+2\alpha_k u.
\]

Регуляризованный метод проекции градиента:
\[
u_{k+1}=P_U\left(u_k-\beta_k(J_k'(u_k)+2\alpha_k u_k)\right).
\]

В задаче нагрева:
\[
J_k'(u_k)\approx\psi_k.
\]
Поэтому:
\[
u_{k+1}=P_U\left(u_k-\beta_k(\psi_k+2\alpha_k u_k)\right).
\]

Ключевая схема:
\[
\begin{gathered}
\text{прямая задача}\Rightarrow\text{сопряженная задача}\Rightarrow J'(u)=\psi\\
\Rightarrow\text{вариационное неравенство}\Rightarrow\text{РМПГ}.
\end{gathered}
\]


\blocktitle{F}{Линейно-квадратичные игры, Berge и Pareto}{Билеты 18--20. Краткая версия для запоминания}

\ticket{18}{Равновесие по Нэшу в ЛК-дифференциальной игре трех лиц}
\official{18. Условия существования и явный вид ситуации равновесия по Нэшу в дифференциальной линейно-квадратичной дифференциальной игре трех лиц.}

Модель:
\[
\dot x=A(t)x+u_1+u_2+u_3,
\qquad x(t_0)=x_0,
\qquad x,u_i\in\R^n.
\]
Игроки используют линейные позиционные стратегии:
\[
u_i(t,x)=Q_i(t)x,
\qquad Q_i(\cdot)\in C^{n\times n}[0,\vartheta].
\]
Функции выигрыша:
\[
J_i=x^{\top}(\vartheta)C_ix(\vartheta)+
\int_{t_0}^{\vartheta}\sum_{j=1}^3 u_j^{\top}D_{ij}u_j\,dt,
\qquad i=1,2,3.
\]
Равновесие по Нэшу:
\[
J_i(U^e,t_0,x_0)\ge J_i(U_i,U_{-i}^e,t_0,x_0)
\quad \forall U_i,
\quad i=1,2,3.
\]

Вводятся функции Беллмана:
\[
V_i(t,x)=x^{\top}\Theta_i(t)x,
\qquad \Theta_i^{\top}=\Theta_i.
\]
Для каждого игрока максимизация идет по его собственному управлению \(u_i\). Условие первого порядка:
\[
\frac{\partial W_i}{\partial u_i}=V_{i,x}+2D_{ii}u_i=0.
\]
Для существования конечного максимума нужно
\[
D_{ii}<0.
\]
Тогда
\[
u_i^e(t,x)=-D_{ii}^{-1}\Theta_i(t)x,
\qquad i=1,2,3.
\]

Матрицы \(\Theta_i\) удовлетворяют связанной системе Риккати:
\[
\begin{aligned}
0={}&\dot\Theta_i+\Theta_iA+A^{\top}\Theta_i-\Theta_iD_{ii}^{-1}\Theta_i \\
&-\sum_{j\ne i}(\Theta_iD_{jj}^{-1}\Theta_j+\Theta_jD_{jj}^{-1}\Theta_i)
+\sum_{j\ne i}\Theta_jD_{jj}^{-1}D_{ij}D_{jj}^{-1}\Theta_j,
\end{aligned}
\]
\[
\Theta_i(\vartheta)=C_i,
\qquad i=1,2,3.
\]
Если эта система имеет продолжимое симметричное решение на \([0,\vartheta]\), то стратегии
\[
U_i^e:\quad u_i^e(t,x)=-D_{ii}^{-1}\Theta_i(t)x
\]
образуют равновесие по Нэшу.

Равновесные выигрыши:
\[
J_i^e(t_0,x_0)=x_0^{\top}\Theta_i(t_0)x_0.
\]
Если хотя бы для одного игрока \(D_{ii}>0\), равновесие по Нэшу отсутствует: максимум по его управлению не существует.

\textbf{Ключевая фраза:} Nash в трехлицевой LQ-игре строится через три функции Беллмана \(V_i=x^{\top}\Theta_i x\); равновесные стратегии имеют вид \(u_i^e=-D_{ii}^{-1}\Theta_i x\), если связанная система Риккати имеет решение.

\newpage
\ticket{19}{Равновесие по Бержу в ЛК-игре двух лиц с малым параметром}
\official{19. Условия существования и явный вид ситуации равновесия по Бержу в дифференциальной линейно-квадратичной игре двух лиц с малым влиянием одного из игроков на скорость изменения фазового вектора.}

Модель:
\[
\dot x=A(t)x+u_1+\varepsilon u_2,
\qquad 0\le\varepsilon\ll1.
\]
Стратегии:
\[
u_i(t,x)=Q_i(t)x,
\qquad i=1,2.
\]
Функции выигрыша:
\[
J_1=x^{\top}(\vartheta)C_1x(\vartheta)+
\int_{t_0}^{\vartheta}(u_1^{\top}D_{11}u_1-u_2^{\top}D_{12}u_2+x^{\top}G_1x)dt,
\]
\[
J_2=x^{\top}(\vartheta)C_2x(\vartheta)+
\int_{t_0}^{\vartheta}(-u_1^{\top}D_{21}u_1+u_2^{\top}D_{22}u_2+x^{\top}G_2x)dt.
\]

Равновесие по Бержу:
\[
J_1(U_1^B,U_2)\le J_1(U^B)\quad \forall U_2,
\]
\[
J_2(U_1,U_2^B)\le J_2(U^B)\quad \forall U_1.
\]
Главное отличие от Нэша:
\[
\begin{aligned}
\text{Нэш: }& J_i\text{ максимизируется по }u_i,\\
\text{Бержу: }& J_i\text{ максимизируется по чужому управлению.}
\end{aligned}
\]
Значит,
\[
J_1\text{ максимизируется по }u_2,
\qquad
J_2\text{ максимизируется по }u_1.
\]

Функции Беллмана:
\[
V_i(t,x)=x^{\top}\Theta_i(t,\varepsilon)x.
\]
Условия первого порядка:
\[
\frac{\partial W_1}{\partial u_2}=\varepsilon V_{1,x}-2D_{12}u_2=0,
\]
\[
\frac{\partial W_2}{\partial u_1}=V_{2,x}-2D_{21}u_1=0.
\]
Отсюда явный вид равновесных стратегий:
\[
u_1^B(t,x)=D_{21}^{-1}\Theta_2(t,\varepsilon)x,
\]
\[
u_2^B(t,x)=\varepsilon D_{12}^{-1}\Theta_1(t,\varepsilon)x.
\]

Условия существования в используемой конвенции:
\[
C_2\le0,
\qquad G_2\le0,
\qquad D_{ij}>0,
\quad i,j=1,2,
\]
и \(\varepsilon\) достаточно мало. Тогда существует равновесие по Бержу--Вайсману, а равновесие по Нэшу в этой постановке не существует.

Выигрыши:
\[
J_i^B(t_0,x_0)=x_0^{\top}\Theta_i(t_0,\varepsilon)x_0,
\qquad i=1,2.
\]
Малый параметр нужен для метода Пуанкаре:
\[
\Theta_i(t,\varepsilon)=\Theta_i^{(0)}(t)+\varepsilon\Theta_i^{(1)}(t)+\varepsilon^2\Theta_i^{(2)}(t)+\cdots .
\]

\textbf{Ключевая фраза:} в равновесии по Бержу выигрыш игрока максимизируется стратегией другого игрока; поэтому \(u_1^B\) зависит от \(\Theta_2\), а \(u_2^B\) зависит от \(\Theta_1\) и содержит множитель \(\varepsilon\).

\newpage
\ticket{20}{Максимум по Парето в двухкритериальной линейно-квадратичной динамической задаче}
\official{20. Максимум по Парето в двухкритериальной линейно-квадратичной динамической задаче.}

Есть два критерия:
\[
f(x,y)\to\max,
\qquad R(x,y)\to\min.
\]
Чтобы оба критерия были на максимум, пишем
\[
G_2=\langle X,Y,\{f(x,y),-R(x,y)\}\rangle .
\]
Здесь \(x\in X\) -- стратегия, \(y\in Y\) -- неопределенность.

Линейно-квадратичный выигрыш:
\[
f(x,y)=x^{\top}Ax+2x^{\top}By+y^{\top}Dy+2a^{\top}x+2d^{\top}y,
\qquad A=A^{\top},
\quad D=D^{\top}.
\]
Риск, или сожаление:
\[
R(x,y)=\max_{z\in X}f(z,y)-f(x,y)\ge0.
\]
Смысл: риск равен потере выигрыша по сравнению с лучшим выбором при уже известном \(y\).

Для фиксированного \(x\) неопределенность \(y^P(x)\) минимальна по Парето, если нельзя одновременно уменьшить выигрыш и увеличить риск. То есть не существует \(y\in Y\), для которого
\[
f(x,y)\le f(x,y^P(x)),
\]
\[
-R(x,y)\le -R(x,y^P(x)),
\]
и хотя бы одно неравенство строгое.

Достаточная конструкция:
\[
f[y]=\max_{x\in X}f(x,y),
\]
\[
y^P(x)\in\arg\min_{y\in Y}\bigl(f(x,y)-\sigma f[y]\bigr),
\qquad \sigma\in(0,1).
\]
Тогда \(y^P(x)\) является минимальной по Парето неопределенностью.

Стратегия \(x^P\) максимальна по Парето, если не существует \(x\in X\), такого что
\[
f(x,y^P(x))\ge f(x^P,y^P(x^P)),
\]
\[
-R(x,y^P(x))\ge -R(x^P,y^P(x^P)),
\]
и хотя бы одно неравенство строгое.

Динамическая версия: вместо \(f\) берутся два функционала
\[
J_1(U,t_0,x_0),
\qquad J_2(U,t_0,x_0).
\]
Pareto максимум: нет другого \(U\), который одновременно улучшает \(J_1\) и \(J_2\).

Положительная скаляризация:
\[
J_\alpha(U)=J_1(U)+\alpha J_2(U),
\qquad \alpha>0.
\]
Если \(U^P\) максимизирует \(J_\alpha\), то \(U^P\) является максимальной по Парето альтернативой.

В LQ-случае после скаляризации:
\[
C(\alpha)=C_1+
\alpha C_2,
\qquad
D_j(\alpha)=D_{1j}+\alpha D_{2j}.
\]
Ищем
\[
V^P(t,x)=x^{\top}\Theta^P(t)x.
\]
Pareto-максимальное управление имеет линейный вид
\[
u_j^P(t,x)=-D_j^{-1}(\alpha)\Theta^P(t)x.
\]
Матрица \(\Theta^P\) находится из уравнения Риккати:
\[
\dot\Theta^P+
\Theta^P A+A^{\top}\Theta^P-
\Theta^P\left(\sum_jD_j^{-1}(\alpha)\right)\Theta^P=0,
\]
\[
\Theta^P(\vartheta)=C(\alpha).
\]

\textbf{Ключевая фраза:} максимум по Парето -- это не один абсолютный максимум, а неулучшаемый компромисс между выигрышем и риском; в LQ-динамике он строится через положительную скаляризацию и Riccati.

\section*{Общая схема блока F}
\addcontentsline{toc}{section}{Общая схема блока F}
\[
\begin{gathered}
\text{LQ игра}\Rightarrow \text{Bellman-функция }V=x^{\top}\Theta x\\
\Rightarrow \text{условие максимума}\Rightarrow \text{Riccati}\\
\Rightarrow \text{линейная позиционная стратегия}.
\end{gathered}
\]
\[
\text{Nash: свой выигрыш по своему управлению.}
\]
\[
\text{Berge: свой выигрыш по чужому управлению.}
\]
\[
\text{Pareto: нельзя улучшить один критерий без ухудшения другого.}
\]


\blocktitle{G}{Риск, гарантированные решения и Cournot}{Билеты 21--23. Краткая версия для запоминания}

\section*{Билет 21. Необходимые и достаточные условия существования стратегий с нулевым риском}
Рассматривается задача при неопределенности
\[
\langle X,Y,f(x,y)\rangle,
\]
где $x\in X$ --- стратегия ЛПР, $y\in Y$ --- неопределенность, $f(x,y)$ --- выигрыш.

Функция риска, или сожаления:
\[
\Phi(x,y)=\max_{z\in X}f(z,y)-f(x,y).
\]
Она всегда неотрицательна:
\[
\Phi(x,y)\ge 0.
\]
Риск равен нулю тогда и только тогда, когда $x$ оптимальна при данном $y$:
\[
\Phi(x,y)=0
\Longleftrightarrow
f(x,y)=\max_{z\in X}f(z,y).
\]

Стратегия $x^0$ имеет нулевой риск, если
\[
\Phi(x^0,y)=0\qquad \forall y\in Y.
\]
То есть
\[
f(x^0,y)=\max_{z\in X}f(z,y)
\qquad \forall y\in Y.
\]

Главное необходимое и достаточное условие:
\[
\exists x^0\in X:\ \Phi(x^0,y)=0\ \forall y
\]
эквивалентно
\[
\bigcap_{y\in Y}\Argmax_{z\in X}f(z,y)\ne\varnothing.
\]

Эквивалентная минимаксная форма:
\[
\min_{x\in X}\max_{y\in Y}\Phi(x,y)=0.
\]
Если используется язык седловой точки, то должна существовать седловая точка нулевого значения:
\[
\max_{y\in Y}\Phi(x^0,y)=\Phi(x^0,y^0)=\min_{x\in X}\Phi(x,y^0)=0.
\]
Правильное минимаксное неравенство:
\[
\max_{y\in Y}\min_{x\in X}\Phi(x,y)
\le
\min_{x\in X}\max_{y\in Y}\Phi(x,y).
\]

Итоговая фраза:
\[
\begin{gathered}
\text{нулевой риск }\Longleftrightarrow\\
\text{одна и та же стратегия оптимальна при всех неопределенностях}.
\end{gathered}
\]

\newpage
\section*{Билет 22. Гарантированное по Парето равновесие в многошаговой дуополии Курно с учетом импорта}
Две фирмы выбирают объемы выпуска, импорт является неопределенностью.

Статическая модель:
\[
q_1,q_2\ge0,
\qquad y\ge0,
\]
\[
\bar q=q_1+q_2+y,
\qquad
p=a-b\bar q,
\quad a>0,\ b>0.
\]
Прибыль фирмы $i$:
\[
\psi_i(q_1,q_2,y)=\bigl[a-b(q_1+q_2+y)\bigr]q_i-(cq_i+d).
\]
Для учета неопределенности вводится гарантированный критерий:
\[
\Phi_i(q_1,q_2,y)=\psi_i(q_1,q_2,y)+\frac{y^2}{2}.
\]

Парето-гарантированное равновесие строится в три шага:
\begin{enumerate}
  \item для каждого $q=(q_1,q_2)$ найти Парето-минимальную неопределенность $y^P(q)$;
  \item подставить $y=y^P(q)$ и найти равновесие Нэша в игре гарантий;
  \item вычислить прибыли и гарантированные выигрыши.
\end{enumerate}

В статической модели:
\[
y^P(q)=\frac{b(q_1+q_2)}{2},
\]
\[
q_1^e=q_2^e=\frac{a-c}{b(3+b)}.
\]
Тогда
\[
y^P(q^e)=\frac{a-c}{3+b},
\]
\[
\psi_i^e=\frac{(a-c)^2}{b(3+b)^2}-d,
\qquad i=1,2.
\]

Многошаговая версия: состояние
\[
q(k)=(q_1(k),q_2(k)).
\]
Двухшаговая модель:
\[
q_1(k+1)=\frac{a-c}{2b}-\frac{q_2(k)}{2}-\frac{z(k)}{2}+u_1(k,q(k)),
\]
\[
q_2(k+1)=\frac{a-c}{2b}-\frac{q_1(k)}{2}-\frac{z(k)}{2}+u_2(k,q(k)),
\qquad k=0,1.
\]
Выигрыши:
\[
J_i(U,Z,q_0)=[q_i(2)]^2+\sum_{k=0}^{1}\left(q_i(k)-2u_i^2(k)+\frac{z^2(k)}{2}\right).
\]

Определение:
\[
(U^e,J_1^e[q_0],J_2^e[q_0])
\]
--- гарантированное по Парето равновесие, если существует $Z^P$, такое что:
\[
Z^P \text{ минимальна по Парето по неопределенности},
\]
\[
U^e \text{ является равновесием Нэша при } Z=Z^P,
\]
\[
J_i^e[q_0]=J_i(U^e,Z^P,q_0).
\]
Алгоритм: обратная индукция
\[
k=2\rightarrow k=1\rightarrow k=0.
\]
Терминальные функции:
\[
V_i^{(2)}(q)=q_i^2.
\]
На каждом предыдущем шаге сначала ищем $z^P(k,q,u)$, затем равновесные управления $u_i^e(k,q)$, затем функции $V_i^{(k)}(q)$. В итоге
\[
J_i^e[q_0]=V_i^{(0)}(q_0).
\]

Главная фраза:
\[
\begin{gathered}
\text{сначала Парето-минимальная неопределенность импорта,}\\
\text{затем Нэш методом обратной индукции.}
\end{gathered}
\]

\newpage
\section*{Билет 23. Гарантированное по выигрышам и рискам равновесие в ЛК-игре двух лиц}
Игра двух лиц при неопределенности:
\[
X_1=X_2=\R^n,
\qquad Y=\R^m.
\]
Выигрыши:
\[
f_1=x_1^\top A_1x_1+2x_1^\top x_2+y^\top D_1y+2a_1^\top x_1+2d_1^\top y,
\]
\[
f_2=x_2^\top A_2x_2-2x_2^\top x_1+y^\top D_2y+2a_2^\top x_2+2d_2^\top y.
\]
Условия:
\[
A_i=A_i^\top<0,
\qquad D_i=D_i^\top>0.
\]
$A_i<0$ нужно для существования максимума по собственной стратегии, $D_i>0$ --- для корректного построения $y^P$.

Лучшие выигрыши:
\[
f_1[x_2,y]=\max_{x_1}f_1(x_1,x_2,y),
\]
\[
f_2[x_1,y]=\max_{x_2}f_2(x_1,x_2,y).
\]
При $A_i<0$:
\[
f_1[x_2,y]=-(x_2+a_1)^\top A_1^{-1}(x_2+a_1)+y^\top D_1y+2d_1^\top y,
\]
\[
f_2[x_1,y]=-(x_1-a_2)^\top A_2^{-1}(x_1-a_2)+y^\top D_2y+2d_2^\top y.
\]

Риски:
\[
R_1=f_1[x_2,y]-f_1(x_1,x_2,y),
\]
\[
R_2=f_2[x_1,y]-f_2(x_1,x_2,y).
\]
Всегда $R_i\ge0$.

Гарантированное по выигрышам и рискам равновесие --- пятерка
\[
(x^e,f_1^P,f_2^P,R_1^P,R_2^P),
\]
для которой существует $y^P$, такое что
\[
f_i^P=f_i(x^e,y^P),
\qquad
R_i^P=R_i(x^e,y^P).
\]
Условия:
\begin{enumerate}
  \item $x^e$ --- равновесие Нэша при фиксированном $y=y^P$;
  \item $y^P$ --- минимальная по Парето неопределенность в четырехкритериальной задаче
  \[
  \{f_1,-R_1,f_2,-R_2\}.
  \]
\end{enumerate}

В данной ЛК-модели
\[
y^P=-(D_1+D_2)^{-1}(d_1+d_2).
\]
Равновесие Нэша при $y=y^P$ находится из системы
\[
A_1x_1^e+x_2^e=-a_1,
\]
\[
-x_1^e+A_2x_2^e=-a_2.
\]
При выбранной конвенции знаков:
\[
x_1^e=(A_2A_1+E_n)^{-1}(a_2-A_2a_1),
\]
\[
x_2^e=-(A_1A_2+E_n)^{-1}(a_1+A_1a_2).
\]
Гарантированные выигрыши и риски:
\[
f_i^P=f_i(x^e,y^P),
\qquad
R_i^P=R_i(x^e,y^P).
\]
Так как $x^e$ --- Nash при $y^P$, то
\[
R_1^P=R_2^P=0.
\]

Главная фраза:
\[
\text{сначала }y^P\text{ по критериям }\{f_1,-R_1,f_2,-R_2\},
\text{ затем Нэш при фиксированном }y^P.
\]


\blocktitle{H}{Экономические модели технологии, ресурсов и климата}{Билеты 24--26. Краткая версия для запоминания}

\section{Билет 24. Математические модели технического прогресса}

\textbf{Официальная формулировка:} Математические модели технического прогресса (экзогенный и эндогенный прогресс). Задачи оптимизации подходы к решению.

Технический прогресс -- это рост выпуска, не сводимый только к росту капитала $K$ и труда $L$. Главное деление:
\[
\text{экзогенный прогресс}\qquad\text{и}\qquad\text{эндогенный прогресс}.
\]

\subsection*{1. Экзогенный технический прогресс}
Технология задается извне:
\[
Y(t)=F(K(t),A(t)L(t)),
\]
\[
\dot A(t)=gA(t),\qquad \dot L(t)=nL(t).
\]
Капитал:
\[
\dot K(t)=I(t)-\delta K(t).
\]
Если $I(t)=sY(t)$, то модель Солоу--Свона:
\[
\dot K=sF(K,AL)-\delta K.
\]
В переменных на единицу эффективного труда
\[
k=\frac{K}{AL}
\]
получаем
\[
\dot k=sf(k)-(\delta+n+g)k.
\]
Смысл: $g$ -- внешний темп технического прогресса.

\subsection*{2. Ramsey с экзогенным прогрессом}
Если выбирается потребление, получаем задачу оптимального управления:
\[
\max_{C(\cdot)}\int_0^\infty e^{-\rho t}U(c(t))L(t)\,dt
\]
при
\[
\dot K=F(K,AL)-C-\delta K.
\]
Current-value Hamiltonian:
\[
\mathcal H=U(c)+\lambda[f(k)-c-(\delta+n+g)k].
\]
Условие максимума:
\[
U'(c)=\lambda.
\]
Сопряженное уравнение:
\[
\dot\lambda=\lambda(\rho+\delta+n+g-f'(k)).
\]
Трансверсальность:
\[
\lim_{t\to\infty}e^{-\rho t}\lambda(t)k(t)=0.
\]

\subsection*{3. Эндогенный технический прогресс}
Технология возникает внутри модели:
\[
Y=F(K,L,A,H).
\]
Здесь $A(t)$ -- знания / технология, $H(t)$ -- человеческий капитал.

Пример модели с человеческим капиталом:
\[
Y=K^\alpha(uH)^{1-\alpha},
\]
\[
\dot K=Y-C-\delta K,
\]
\[
\dot H=B(1-u)H.
\]
Управления:
\[
C(t),\quad u(t).
\]
Смысл: $u$ -- доля человеческого капитала в производстве, $1-u$ -- доля в образовании.

\subsection*{4. R\&D / знания}
Можно писать:
\[
\dot A=\Phi(A,L_A,A_F),
\]
например
\[
\dot A=\beta L_AA
\]
или
\[
\dot A=\beta L_A(A_F-A).
\]
Здесь $L_A$ -- труд в секторе НИОКР, $A_F$ -- внешний технологический уровень.

\subsection*{5. Подходы к решению}
\begin{enumerate}
\item Принцип максимума Понтрягина:
\[
H=F_0(x,u)+\psi^\top f(x,u),
\]
\[
\dot\psi=\rho\psi-H_x,
\]
\[
H(x^*,u^*,\psi)=\max_u H(x^*,u,\psi).
\]
\item Метод Беллмана:
\[
\rho V(x)=\max_u\{F_0(x,u)+V_xf(x,u)\}.
\]
\item Анализ steady state / balanced growth path.
\end{enumerate}

\textbf{Главная фраза:}
\begin{center}
\fbox{\begin{minipage}{0.92\textwidth}
В экзогенных моделях $A(t)$ задана извне; в эндогенных моделях $A(t)$ или $H(t)$ являются состояниями и зависят от инвестиций в знания, образование и НИОКР.
\end{minipage}}
\end{center}

\section{Билет 25. Оптимальная добыча невозобновляемых ресурсов}

\textbf{Официальная формулировка:} Математические модели оптимальной добычи невозобновляемых ресурсов. Задачи оптимизации подходы к решению.

Невозобновляемый ресурс не восстанавливается на рассматриваемом горизонте.

Фазовая переменная:
\[
R(t)\ge 0
\]
-- запас ресурса.

Управление:
\[
E(t)
\]
-- интенсивность добычи.

Динамика:
\[
\dot R(t)=-E(t),\qquad R(0)=R_0>0.
\]
Ограничение:
\[
0\le E(t)\le E_{\max}.
\]
Функционал прибыли:
\[
J(E,T)=\int_0^T e^{-rt}[p(t)E(t)-C(E(t),R(t))]dt\to\max.
\]
Возможные условия:
\[
R(T)=0
\]
при полном исчерпании; $T$ может быть заданным или свободным.

\subsection*{ПМП}
Гамильтониан:
\[
H=e^{-rt}[p(t)E-C(E,R)]-\psi E.
\]
Сопряженное уравнение:
\[
\dot\psi=-H_R=e^{-rt}C_R(E,R).
\]
Условие максимума:
\[
E^*(t)\in\argmax_{0\le E\le E_{\max}}\{e^{-rt}[p(t)E-C(E,R^*)]-\psi E\}.
\]
Если максимум внутренний:
\[
e^{-rt}[p(t)-C_E(E^*,R^*)]=\psi(t).
\]
Экономический смысл: $\psi(t)$ -- теневая цена единицы ресурса в недрах.

\subsection*{Bang-bang случай}
Если $C=0$, то
\[
H=(e^{-rt}p(t)-\psi(t))E.
\]
Поэтому
\[
e^{-rt}p(t)>\psi(t)\Rightarrow E^*(t)=E_{\max},
\]
\[
e^{-rt}p(t)<\psi(t)\Rightarrow E^*(t)=0.
\]
При равенстве возможен особый режим.

\subsection*{Правило Хотеллинга}
Чистая цена ресурса:
\[
\pi(t)=p(t)-c(t).
\]
Правило Хотеллинга:
\[
\frac{\dot\pi(t)}{\pi(t)}=r.
\]
Если $c(t)=0$, то
\[
\frac{\dot p(t)}{p(t)}=r,
\]
\[
p(t)=p_0e^{rt}.
\]
Смысл: ресурс в недрах должен приносить ту же доходность, что и финансовый актив.

\subsection*{Модель Дасгупты--Хила}
Ресурс входит в производство:
\[
Y=F(K,E).
\]
Баланс:
\[
Y=C+I.
\]
Капитал:
\[
\dot K=I-\delta K.
\]
Ресурс:
\[
\dot R=-E.
\]
Задача:
\[
\max\int_0^\infty e^{-\rho t}U(C(t))dt.
\]
Вывод: оптимальная политика балансирует текущее потребление, накопление капитала и сохранение ресурса.

\textbf{Главная фраза:}
\begin{center}
\fbox{\begin{minipage}{0.92\textwidth}
Оптимальная добыча определяется сравнением текущей предельной ренты от добычи с теневой ценой ресурса; в классической модели Хотеллинга чистая цена ресурса растет с темпом $r$.
\end{minipage}}
\end{center}

\section{Билет 26. Интегрированные модели климата и экономики. DICE}

\textbf{Официальная формулировка:} Интегрированные модели климата и экономики. Модель DICE. Задачи оптимизации подходы к решению.

DICE:
\[
\text{Dynamic Integrated model of Climate and the Economy}.
\]
Это integrated assessment model, соединяющая
\[
\begin{gathered}
\text{экономический рост}\to\text{выбросы}\to\text{углеродный цикл}\\
\to\text{температура}\to\text{ущерб}\to\text{общественное благосостояние}.
\end{gathered}
\]

\subsection*{1. Целевая функция}
DICE максимизирует общественное благосостояние:
\[
W=\sum_{t=1}^{T}R(t)U(c(t),L(t))\to\max.
\]
Обычно
\[
U(c,L)=L\frac{c^{1-\alpha}}{1-\alpha},
\]
\[
c(t)=\frac{C(t)}{L(t)}.
\]
Дисконтирование:
\[
R(t)=(1+\rho)^{-t}.
\]

\subsection*{2. Экономический блок}
Валовой выпуск:
\[
Y_{\mathrm{gross}}(t)=A(t)K^\gamma(t)L^{1-\gamma}(t).
\]
Капитал:
\[
K(t+1)=(1-\delta)^hK(t)+hI(t).
\]
Баланс:
\[
Y(t)=C(t)+I(t).
\]
Управление:
\[
s(t)=\frac{I(t)}{Y(t)}.
\]

\subsection*{3. Ущерб и издержки снижения выбросов}
Климатический ущерб:
\[
\Omega(t)=\psi_1T_{\mathrm{AT}}(t)+\psi_2T_{\mathrm{AT}}^2(t).
\]
Издержки снижения выбросов:
\[
\Lambda(t)=\theta_1(t)\mu(t)^{\theta_2}.
\]
Чистый выпуск:
\[
Y(t)=\frac{1-\Lambda(t)}{1+\Omega(t)}Y_{\mathrm{gross}}(t).
\]

\subsection*{4. Выбросы}
Промышленные выбросы:
\[
E_{\mathrm{ind}}(t)=\sigma(t)(1-\mu(t))Y_{\mathrm{gross}}(t).
\]
Полные выбросы:
\[
E(t)=E_{\mathrm{ind}}(t)+E_{\mathrm{land}}(t).
\]
Управление:
\[
\mu(t)\in[0,1]
\]
-- emissions control rate.

\subsection*{5. Углеродный цикл}
Три резервуара:
\[
M_{\mathrm{AT}},\quad M_{\mathrm{UP}},\quad M_{\mathrm{LO}}.
\]
Схема:
\[
M(t+1)=\Phi M(t)+BE(t).
\]
Смысл: выбросы попадают в атмосферу, затем перераспределяются между атмосферой, верхним океаном / биосферой и глубоким океаном.

\subsection*{6. Радиационное воздействие}
\[
F(t)=\eta\log_2\left(\frac{M_{\mathrm{AT}}(t)}{M_{\mathrm{AT}}^{1750}}\right)+F_{\mathrm{EX}}(t).
\]

\subsection*{7. Температурный блок}
Температуры:
\[
T_{\mathrm{AT}}(t)
\]
-- атмосфера / поверхность,
\[
T_{\mathrm{LO}}(t)
\]
-- глубокий океан.

Типично:
\[
T_{\mathrm{AT}}(t+1)=T_{\mathrm{AT}}(t)+\xi_1[F(t)-\xi_2T_{\mathrm{AT}}(t)-\xi_3(T_{\mathrm{AT}}(t)-T_{\mathrm{LO}}(t))],
\]
\[
T_{\mathrm{LO}}(t+1)=T_{\mathrm{LO}}(t)+\xi_4(T_{\mathrm{AT}}(t)-T_{\mathrm{LO}}(t)).
\]

\subsection*{8. Оптимизационная задача}
Управления:
\[
s(t),\quad \mu(t).
\]
Задача:
\[
\max_{\{s(t),\mu(t)\}}\sum_t R(t)L(t)\frac{c(t)^{1-\alpha}}{1-\alpha}
\]
при экономических и климатических ограничениях.

\subsection*{9. Подходы к решению}
DICE решается как дискретная задача нелинейного программирования:
\[
\max W
\]
при системе нелинейных уравнений и ограничений. На практике используются GAMS / NLP solver / CONOPT / Excel Solver. Bellman и ПМП можно использовать как теоретическую интерпретацию, но основная реализация DICE -- прямое нелинейное программирование.

\textbf{Главная фраза:}
\begin{center}
\fbox{\begin{minipage}{0.92\textwidth}
DICE выбирает траектории сбережений $s(t)$ и сокращения выбросов $\mu(t)$, максимизируя дисконтированное благосостояние и связывая экономический рост с углеродным циклом, температурой и климатическим ущербом.
\end{minipage}}
\end{center}


\blocktitle{I}{Дифференциальные игры и конкурентная реклама}{Билеты 27--29. Краткая версия для запоминания}

\section{Билет 27. Неантагонистические дифференциальные игры. Nash в программных и позиционных стратегиях}

Игра $N$ лиц:
\[
 \dot x=f(t,x,u_1,\ldots,u_N),\qquad x(t_0)=x_0.
\]
Выигрыш игрока $i$:
\[
 J_i=\int_{t_0}^{T}F_i(t,x,u_1,\ldots,u_N)\,dt+q_i(x(T)).
\]
Игра неантагонистическая: выигрыши игроков не обязаны быть противоположными, то есть это не нулевая сумма.

\textbf{Равновесие по Нэшу:}
\[
 u^*=(u_1^*,\ldots,u_N^*)
\]
таково, что
\[
 J_i(u_i,u_{-i}^*)\le J_i(u_i^*,u_{-i}^*)\qquad \forall u_i,\quad i=1,\ldots,N.
\]

\subsection*{1. Программные стратегии}

\[
 u_i=u_i(t).
\]
Они зависят от времени, но не от текущего состояния.

Гамильтониан игрока $i$:
\[
 H_i=F_i+\langle\psi_i,f\rangle.
\]
У каждого игрока своя сопряженная переменная $\psi_i$.

Необходимые условия PMP:
\[
 \dot x^*=f(t,x^*,u_1^*,\ldots,u_N^*),
\]
\[
 \dot\psi_i=-H_{i,x},
\]
\[
 \psi_i(T)=q_{i,x}(x^*(T)),
\]
\[
 u_i^*(t)\in\argmax_{v_i\in U_i}H_i(t,x^*,u_1^*,\ldots,u_{i-1}^*,v_i,u_{i+1}^*,\ldots,u_N^*,\psi_i).
\]

\subsection*{2. Позиционные стратегии}

\[
 u_i=\varphi_i(t,x).
\]
Они зависят от текущего состояния. Состоятельное позиционное равновесие описывается системой HJB.

Функции Беллмана:
\[
 V_i(t,x),\qquad V_i(T,x)=q_i(x).
\]
Система:
\[
 -V_{i,t}(t,x)=\max_{v_i\in U_i}\{F_i(t,x,\varphi_{-i}^*,v_i)+V_{i,x}(t,x)f(t,x,\varphi_{-i}^*,v_i)\}.
\]
Равновесная обратная связь:
\[
 \varphi_i^*(t,x)\in\argmax_{v_i\in U_i}\{F_i(t,x,\varphi_{-i}^*,v_i)+V_{i,x}f(t,x,\varphi_{-i}^*,v_i)\}.
\]

\textbf{Главная фраза:}
\begin{quote}
Program Nash: PMP для каждого игрока. Feedback Nash: система HJB.
\end{quote}

\newpage
\section{Билет 28. Дифференциальные игры с бесконечной продолжительностью}

Игра $N$ лиц на $[t_0,+\infty)$:
\[
 \dot x=f(t,x,u_1,\ldots,u_N),\qquad x(t_0)=x_0.
\]
Выигрыш игрока $i$:
\[
 J_i=\int_{t_0}^{+\infty}e^{-r(t-t_0)}F_i(t,x,u_1,\ldots,u_N)\,dt.
\]
Нужно обеспечить сходимость $J_i$; одного условия $r>0$ вообще недостаточно.

\subsection*{1. Программные стратегии}

\[
 u_i=u_i(t).
\]
Nash equilibrium:
\[
 J_i(u_i,u_{-i}^*)\le J_i(u_i^*,u_{-i}^*)\qquad \forall u_i,\quad i=1,\ldots,N.
\]

Present-value Hamiltonian:
\[
 \mathcal H_i=e^{-r(t-t_0)}F_i+\langle\psi_i,f\rangle.
\]
Current-value переменная:
\[
 p_i(t)=e^{r(t-t_0)}\psi_i(t),\qquad \psi_i(t)=e^{-r(t-t_0)}p_i(t).
\]
Current-value Hamiltonian:
\[
 H_i=F_i+\langle p_i,f\rangle.
\]
Необходимые условия:
\[
 u_i^*(t)\in\argmax_{v_i\in U_i}H_i(t,x^*,u_{-i}^*,v_i,p_i),
\]
\[
 \dot p_i=rp_i-\frac{\partial H_i}{\partial x}.
\]
На бесконечности требуется условие трансверсальности, например
\[
 \lim_{t\to\infty}\langle\psi_i(t),x^*(t)\rangle=0,
\]
или другая подходящая форма в конкретной модели.

\subsection*{2. Позиционные стратегии}

\[
 u_i=\varphi_i(t,x),
\]
а в автономной задаче обычно
\[
 u_i=\varphi_i(x).
\]
Состоятельное позиционное равновесие описывается стационарной системой Беллмана:
\[
 rV_i(x)=\max_{v_i\in U_i}\{F_i(x,\varphi_{-i}^*(x),v_i)+V_{i,x}(x)f(x,\varphi_{-i}^*(x),v_i)\}.
\]
Равновесная стратегия:
\[
 \varphi_i^*(x)\in\argmax_{v_i\in U_i}\{F_i(x,\varphi_{-i}^*(x),v_i)+V_{i,x}(x)f(x,\varphi_{-i}^*(x),v_i)\}.
\]

\textbf{Главная фраза:}
\begin{quote}
На бесконечном горизонте программный Nash задается PMP с current-value переменными и условием трансверсальности; позиционный Nash задается стационарной системой Беллмана: $rV_i=\max\{F_i+V_{i,x}f\}$.
\end{quote}

\newpage
\section{Билет 29. Модель конкурентной рекламы с двумя участниками}

Две фирмы конкурируют на рынке фиксированного размера.
\[
 x(t)\in[0,1]
\]
-- доля рынка первой фирмы,
\[
 1-x(t)
\]
-- доля второй фирмы.

Управления:
\[
 u_1(t),u_2(t)\ge0
\]
-- расходы на рекламу.

Динамика:
\[
 \dot x=u_1\sqrt{1-x}-u_2\sqrt{x},\qquad x(0)=x_0.
\]
Реклама первой фирмы увеличивает $x$, реклама второй уменьшает $x$.

\subsection*{Функционалы}

\[
 J_1=\int_0^T\left(q_1x-\frac{c_1}{2}u_1^2\right)e^{-rt}\,dt+e^{-rT}S_1x(T).
\]
\[
 J_2=\int_0^T\left(q_2(1-x)-\frac{c_2}{2}u_2^2\right)e^{-rt}\,dt+e^{-rT}S_2(1-x(T)).
\]

\subsection*{1. Программное равновесие}

Hamiltonian первой фирмы:
\[
 H_1=e^{-rt}\left(q_1x-\frac{c_1}{2}u_1^2\right)+\Lambda_1(u_1\sqrt{1-x}-u_2\sqrt{x}).
\]
Hamiltonian второй фирмы:
\[
 H_2=e^{-rt}\left(q_2(1-x)-\frac{c_2}{2}u_2^2\right)+\Lambda_2(u_1\sqrt{1-x}-u_2\sqrt{x}).
\]

Условия максимума:
\[
 u_1^*=\frac{\Lambda_1}{c_1}e^{rt}\sqrt{1-x^*},
\]
\[
 u_2^*=-\frac{\Lambda_2}{c_2}e^{rt}\sqrt{x^*}.
\]
При ограничениях $u_i\ge0$ нужно брать положительную часть этих выражений.

Сопряженные уравнения:
\[
 \dot\Lambda_1=-q_1e^{-rt}+\Lambda_1\left(\frac{u_1^*}{2\sqrt{1-x^*}}+\frac{u_2^*}{2\sqrt{x^*}}\right),
\]
\[
 \dot\Lambda_2=q_2e^{-rt}+\Lambda_2\left(\frac{u_1^*}{2\sqrt{1-x^*}}+\frac{u_2^*}{2\sqrt{x^*}}\right).
\]
Терминальные условия:
\[
 \Lambda_1(T)=e^{-rT}S_1,
\]
\[
 \Lambda_2(T)=-e^{-rT}S_2.
\]

\subsection*{2. Позиционное равновесие}

Ищем функции Беллмана:
\[
 V_1(t,x)=e^{-rt}(A_1(t)x+B_1(t)),
\]
\[
 V_2(t,x)=e^{-rt}(A_2(t)(1-x)+B_2(t)).
\]
Тогда
\[
 V_{1x}=e^{-rt}A_1(t),\qquad V_{2x}=-e^{-rt}A_2(t).
\]

Равновесные обратные связи:
\[
 \varphi_1^*(t,x)=\frac{A_1(t)}{c_1}\sqrt{1-x},
\]
\[
 \varphi_2^*(t,x)=\frac{A_2(t)}{c_2}\sqrt{x}.
\]

Коэффициенты определяются системой:
\[
 \dot A_1=rA_1-q_1+\frac{A_1^2}{2c_1}+\frac{A_1A_2}{c_2},
\]
\[
 \dot B_1=rB_1-\frac{A_1^2}{2c_1},
\]
\[
 \dot A_2=rA_2-q_2+\frac{A_1A_2}{c_1}+\frac{A_2^2}{2c_2},
\]
\[
 \dot B_2=rB_2-\frac{A_2^2}{2c_2}.
\]
Терминальные условия:
\[
 A_1(T)=S_1,\quad B_1(T)=0,
\]
\[
 A_2(T)=S_2,\quad B_2(T)=0.
\]

Равновесная динамика:
\[
 \dot x=\frac{A_1(t)}{c_1}(1-x)-\frac{A_2(t)}{c_2}x.
\]

\textbf{Главная фраза:}
\begin{quote}
В модели конкурентной рекламы программное равновесие строится по PMP, а позиционное -- через линейные по доле рынка функции Беллмана.
\end{quote}


\end{document}
